\subsection{Harness-Level Continuation Ablations}
\label{sec:ablation-experiment-setting}
\label{sec:appendix-ablation-setting}

Long-running agent evaluation depends not only on model capability, but also on how the harness keeps the run alive and carries useful state across many hours. In a 12-hour task, an agent may stop its running unexpectedly due to idleness and need to resume. These continuation mechanisms are therefore part of the practical measurement setup: a weak scaffold can make a capable model stop working on the task, while a stronger scaffold may help the same model keep active and improving. As a supplementary harness-level ablation, we evaluate two such mechanisms, /goal mode and the Ralph loop, under a fixed 12-hour budget with GPT-5.5 and GPT-5.4.

Base uses the standard harness: one continuing agent session, a stop hook to prevent voluntary early exit, and auto-resume for abnormal exits. In /goal mode~\cite{openai2026codexgoals}, the harness prompts the agent to create a task-level goal at the beginning, keep it active during the run, and mark it complete only when the task is validated. The Ralph loop follows the file-backed fresh-context pattern~\cite{huntley2026ralphloop}: each loop starts a new agent invocation on the same workspace, asks it to read and update progress.md, then appends judge feedback to that file before the next loop. It uses up to 100 loops, a 7200-second per-loop cap, and the same 12-hour overall budget. Each cell schedules three runs and reports the mean score over valid runs.

\providecommand{\NA}{\textemdash}
\providecommand{\best}[1]{}
\renewcommand{\best}[1]{\begingroup\bfseries\boldmath #1\endgroup}
\providecommand{\second}[1]{}
\renewcommand{\second}[1]{\underline{#1}}
\providecommand{\sci}[2]{\ensuremath{#1\times 10^{#2}}}
\begin{table*}[t!]
\centering
\footnotesize
\renewcommand{\arraystretch}{1.08}
\setlength{\tabcolsep}{2.5pt}
\resizebox{\textwidth}{!}{%
\begin{tabular}{@{}>{\raggedright\arraybackslash}p{0.34\linewidth}*{6}{>{\centering\arraybackslash}p{0.10\linewidth}}@{}}
\BenchTopRule
\textbf{Task} & \multicolumn{3}{c}{\textbf{GPT-5.5}} & \multicolumn{3}{c}{\textbf{GPT-5.4}} \\
\BenchCMidRule{2-4} \BenchCMidRule{5-7}
 & \textbf{Base} & \textbf{w/ Goal} & \textbf{w/ Ralph} & \textbf{Base} & \textbf{w/ Goal} & \textbf{w/ Ralph} \\
\BenchMidRule
\texttt{portfolio\_\allowbreak{}risk\_\allowbreak{}calibration} & \ensuremath{\mathrm{25.0}} & \second{\ensuremath{\mathrm{27.5}}} & \best{\ensuremath{\mathrm{34.3}}} & \ensuremath{\mathrm{10.7}} & \best{\ensuremath{\mathrm{19.8}}} & \second{\ensuremath{\mathrm{12.2}}} \\
\texttt{storyboard\_\allowbreak{}ad\_\allowbreak{}copywriting} & \ensuremath{\mathrm{77.0}} & \second{\ensuremath{\mathrm{88.3}}} & \best{\ensuremath{\mathrm{97.3}}} & \ensuremath{\mathrm{65.7}} & \second{\ensuremath{\mathrm{78.0}}} & \best{\ensuremath{\mathrm{83.3}}} \\
\texttt{arc\_\allowbreak{}compiler\_\allowbreak{}runtime} & \best{\ensuremath{\mathrm{72.4}}} & \second{\ensuremath{\mathrm{70.6}}} & \ensuremath{\mathrm{52.6}} & \second{\ensuremath{\mathrm{50.0}}} & \ensuremath{\mathrm{47.2}} & \best{\ensuremath{\mathrm{51.2}}} \\
\texttt{battery\_\allowbreak{}soh\_\allowbreak{}rul\_\allowbreak{}anomaly} & \second{\ensuremath{\mathrm{30.2}}} & \ensuremath{\mathrm{23.8}} & \best{\ensuremath{\mathrm{48.8}}} & \best{\ensuremath{\mathrm{14.7}}} & \second{\ensuremath{\mathrm{14.6}}} & \ensuremath{\mathrm{13.4}} \\
\texttt{borden\_\allowbreak{}source\_\allowbreak{}inversion} & \second{\ensuremath{\mathrm{38.5}}} & \ensuremath{\mathrm{32.8}} & \best{\ensuremath{\mathrm{62.2}}} & \ensuremath{\mathrm{8.0}} & \second{\ensuremath{\mathrm{17.3}}} & \best{\ensuremath{\mathrm{23.2}}} \\
\texttt{capecod\_\allowbreak{}plume\_\allowbreak{}reconstruction} & \second{\ensuremath{\mathrm{16.4}}} & \ensuremath{\mathrm{16.0}} & \best{\ensuremath{\mathrm{17.3}}} & \ensuremath{\mathrm{12.6}} & \second{\ensuremath{\mathrm{13.3}}} & \best{\ensuremath{\mathrm{13.7}}} \\
\texttt{combinatorial\_\allowbreak{}games\_\allowbreak{}formalization} & \ensuremath{\mathrm{38.2}} & \best{\ensuremath{\mathrm{45.3}}} & \second{\ensuremath{\mathrm{39.8}}} & \second{\ensuremath{\mathrm{17.8}}} & \ensuremath{\mathrm{13.0}} & \best{\ensuremath{\mathrm{20.2}}} \\
\texttt{ffmpeg\_\allowbreak{}swscale\_\allowbreak{}reimplementation} & \ensuremath{\mathrm{15.3}} & \second{\ensuremath{\mathrm{16.4}}} & \best{\ensuremath{\mathrm{25.2}}} & \ensuremath{\mathrm{13.9}} & \best{\ensuremath{\mathrm{28.3}}} & \second{\ensuremath{\mathrm{21.8}}} \\
\texttt{flt\_\allowbreak{}regular\_\allowbreak{}formalization} & \best{\ensuremath{\mathrm{75.1}}} & \second{\ensuremath{\mathrm{66.7}}} & \ensuremath{\mathrm{58.6}} & \best{\ensuremath{\mathrm{48.3}}} & \second{\ensuremath{\mathrm{43.7}}} & \second{\ensuremath{\mathrm{43.7}}} \\
\texttt{git\_\allowbreak{}rewrite\_\allowbreak{}in\_\allowbreak{}zig} & \ensuremath{\mathrm{18.4}} & \second{\ensuremath{\mathrm{20.5}}} & \best{\ensuremath{\mathrm{22.0}}} & \ensuremath{\mathrm{15.4}} & \best{\ensuremath{\mathrm{18.4}}} & \second{\ensuremath{\mathrm{15.8}}} \\
\texttt{tryst\_\allowbreak{}text\_\allowbreak{}adventure} & \second{\ensuremath{\mathrm{55.7}}} & \best{\ensuremath{\mathrm{56.2}}} & \ensuremath{\mathrm{40.2}} & \best{\ensuremath{\mathrm{44.3}}} & \ensuremath{\mathrm{32.9}} & \second{\ensuremath{\mathrm{42.4}}} \\
\texttt{openttd\_\allowbreak{}transport\_\allowbreak{}ai} & \ensuremath{\mathrm{28.1}} & \best{\ensuremath{\mathrm{39.8}}} & \second{\ensuremath{\mathrm{30.6}}} & \best{\ensuremath{\mathrm{11.9}}} & \ensuremath{\mathrm{1.2}} & \second{\ensuremath{\mathrm{8.1}}} \\
\texttt{pocketbase\_\allowbreak{}backend\_\allowbreak{}architecture} & \best{\ensuremath{\mathrm{62.5}}} & \best{\ensuremath{\mathrm{62.5}}} & \second{\ensuremath{\mathrm{41.7}}} & \second{\ensuremath{\mathrm{20.8}}} & \best{\ensuremath{\mathrm{54.2}}} & \ensuremath{\mathrm{0.0}} \\
\texttt{symbolic\_\allowbreak{}integration\_\allowbreak{}engine} & \best{\ensuremath{\mathrm{44.0}}} & \ensuremath{\mathrm{36.6}} & \second{\ensuremath{\mathrm{37.7}}} & \ensuremath{\mathrm{30.9}} & \best{\ensuremath{\mathrm{63.7}}} & \second{\ensuremath{\mathrm{37.2}}} \\
\BenchMidRule
\textbf{Avg.} & \ensuremath{\mathrm{42.6}} & \second{\ensuremath{\mathrm{43.1}}} & \best{\ensuremath{\mathrm{43.4}}} & \ensuremath{\mathrm{26.1}} & \best{\ensuremath{\mathrm{31.8}}} & \second{\ensuremath{\mathrm{27.6}}} \\
\BenchBottomRule
\end{tabular}%
}
\caption{Harness-level continuation ablation of /goal mode and the Ralph loop. GPT-5.5 and GPT-5.4 are evaluated with Base, Goal, and Ralph settings under the same 12-hour budget; each value is the mean score over valid runs, and the Avg. row averages the displayed task rows. Incomplete cells are detailed in Appendix~\ref{sec:appendix-ablation-setting}. Bold marks the best setting and underlining marks the second-best setting.}
\label{tab:feature_ablation_goal_ralph}
\end{table*}


As shown in Table~\ref{tab:feature_ablation_goal_ralph}, the goal mode and the Ralph loop often outperform the base setting, suggesting that long-horizon agents benefit from harness support that preserves and updates task state across extended runs. In the displayed-task average, GPT-5.5 improves from 42.6 in Base to 43.1 with Goal and 43.4 with Ralph, while GPT-5.4 improves from 26.1 in Base to 31.8 with Goal and 27.6 with Ralph. The gains are not uniform across all tasks and models, so we treat this as an appendix-level harness diagnostic rather than a main result.

Most cells use all three scheduled runs. A few GPT-5.4 cells have fewer valid runs because of intermittent API or network instability: one-run cells are w/ Goal for \texttt{storyboard\_\allowbreak{}ad\_\allowbreak{}copywriting} and \texttt{flt\_\allowbreak{}regular\_\allowbreak{}formalization}, plus w/ Ralph for \texttt{symbolic\_integration\_engine}; two-run cells are \texttt{battery\_\allowbreak{}soh\_\allowbreak{}rul\_\allowbreak{}anomaly} (w/ Goal and w/ Ralph), \texttt{capecod\_\allowbreak{}plume\_\allowbreak{}reconstruction} (w/ Goal), \texttt{combinatorial\_\allowbreak{}games\_\allowbreak{}formalization} (w/ Ralph), \texttt{ffmpeg\_\allowbreak{}swscale\_\allowbreak{}reimplementation} (w/ Goal), \texttt{flt\_\allowbreak{}regular\_\allowbreak{}formalization} (w/ Ralph), \texttt{git\_\allowbreak{}rewrite\_\allowbreak{}in\_\allowbreak{}zig} (w/ Goal), and \texttt{symbolic\_integration\_engine} (w/ Goal). All reported GPT-5.5 cells use the full three runs.
